\documentclass {article}

\usepackage [utf8]{inputenc}
\usepackage {graphicx}
\usepackage {listings}
\usepackage [spanish]{babel}
\usepackage {color}

\lstdefinestyle{mystyle}{ 
    backgroundcolor=\color{white}, 
    commentstyle=\color{green}, 
    keywordstyle=\color{blue}, 
    basicstyle=\footnotesize\ttfamily,
    numbers=left,
    numberstyle=\tiny,
    frame=single,
    breaklines=true
}

\lstset{style=mystyle}

\author{Luis Fernando Torres Hernández}
\title{Practica 2}
\date{\today}
\def\description{Cifrado Cesar y listas ligadas simples.}

\makeatletter 
\def\@maketitle{ 
    \raggedright
    \includegraphics[scale=0.5]{logo.png}\\[8ex] 
    \vspace{1.0cm}
\begin{center} {
    \Huge 
    \bfseries
    \sffamily \@title }\\[4ex] {\Large \@author}\\[4ex] 
    \@date\\[8ex]
    \vspace{1.0cm}
    \Large\description
    \newpage
\end{center}} 
\makeatother

\begin{document}
\maketitle


\section {Introducción}
El cifrado Cesar es uno de los primeros métodos para ocultar 
información que se conocen, su funcionamiento es sencillo, 
consiste en un desplazamiento circular del alfabeto. Dadas las
capacidades de computo actuales este cifrado es débil a un 
ataque de fuerza bruta.

Las Listas ligadas simples son una estructura de datos que permite
almacenar datos de forma dinámica, permite insertar elementos 
sin un limite mas que el de la memoria disponible para nuestro 
programa.

\section {Desarrollo}
Esta practica se desarrollo en dos partes, la primera consiste en
un programa escrito en C que implementa el cifrado Cesar y la segunda
parte consiste en un programa que implementa una lista ligada simple
y algunos métodos para manipularla.

\subsection{Parte 1}

El programa desarrollado para el cifrado Cesar tiene 2 argumentos 
obligatorios, el primero es una opción que determina el comportamiento 
del programa \texttt{-c} para cifrar y \texttt{-d} para descifrar,
el segundo argumento es un archivo que contiene el texto que nuestro 
programa va a manipular, el tercer argumento es opcional y es un número 
que indica el desplazamiento que se aplicara al alfabeto, la 
Figura \ref{fig:argumentos} muestra como se manejan estos argumentos
dentro del código.

\begin{figure}[htbp]
    \centering
    \begin{lstlisting}[language=C]
if(argc < 3){
    print_ayuda(argv[0]);
    return 1;
}    

if(argv[1][0] != '-' || (argv[1][1] != 'c' && argv[1][1] != 'd')){
    print_ayuda(argv[0]);
    return 1;
}
archivo = argv[2];

if(argc >= 4){
    llave = atoi(argv[3]) % 26;
}
    
    \end{lstlisting}
    \caption{Manejo de argumentos en archivo p1.c}
    \label{fig:argumentos}
\end{figure}

El cifrado y descifrado son exactamente iguales, la única 
diferencia esta en el desplazamiento que se aplica. Dado este 
comportamiento el programa establece la llave del cifrado en 
\texttt{llave = 26 - llave} cuando se desea descifrar.

La Figura \ref{fig:cifrado} muestra la función que realiza el
cifrado y descifrado, esta función recibe un archivo previamente
abierto y un desplazamiento. escribe el resultado en la salida 
estándar.

\begin{figure}[htbp]
    \centering
    \begin{lstlisting}[language=C]
void cesar(int llave, FILE *entrada){
    char a, c;
    a = 0;
    // Mientras el caracter actual no sea EOF 
    while(!feof(entrada)){
        a = 0;
        // Se lee un caracter del archivo de entrada
        char c = fgetc(entrada);
        if(c == EOF) break; 

        // El caracter leido es una mayuscula
        if(c >= 'A' && c <= 'Z'){
            a = 'A';
        } else if(c >= 'a' && c <= 'z'){
            a = 'a';
        }
        if(a != 0){
            c = (c - a + llave) % 26 + a;
        }
        printf("%c", c);
    }
}
        
    \end{lstlisting}
    \caption{Función de cifrado y descifrado}
    \label{fig:cifrado}
\end{figure}

\subsection{Parte 2}
La segunda parte de la práctica consiste realizar una 
implementación de una lista ligada simple. La implementación 
implica la definición de dos estructuras (\texttt{struct}), 
una para los nodos que conforman la lista y otra para llevar 
el control del primer elemento, estas estructuras se pueden 
ver en la Figura \ref{fig:structs}. Todos los métodos que manipulan
la lista se implementaron mediante funciones que reciben un 
apuntador a una lista como primer argumento, esto permite separar
la lógica de cada método del código de la aplicación principal.

\begin{figure}
    \centering
    \begin{lstlisting}[language=C]
typedef struct nodo {
    int dato;
    struct nodo *siguiente;
}nodo;

typedef struct {
    nodo *inicio;
}lista;
    \end{lstlisting}
    \caption{Definición de estructuras para la lista ligada simple}
    \label{fig:structs}
\end{figure}

Los métodos implementados son los siguientes:
\begin{itemize}
    \item \texttt{insertar(lista *l, int dato)}: 
    Inserta un elemento en la lista manteniendo el orden 
    ascendente de los elementos.

    \item \texttt{desplegar(lista *l)}: 
    Muestra los elementos de la lista en salida estándar.

    \item \texttt{consultar(lista *l, int index)}:
    Regresa un apuntador al elemento en la posición \texttt{index} 
    de la lista, regresa \texttt{NULL} si la posición no es válida.

    \item \texttt{eliminar(lista *l, int index)}: Quita el nodo en la 
    posición \texttt{index} de la lista, si la posición no es válida
    no hace nada y regresa \texttt{NULL}, si la posición es válida 
    regresa un apuntador al nodo eliminado.

    \item \texttt{buscar(lista *l, int dato)}: 
    Busca un elemento en la lista y regresa el índice donde se encuentra,
    regresa -1 si el elemento no se encuentra en la lista.
    
    \item \texttt{buscar\_recursivo(nodo *n)}:
    Busca un elemento en la lista de forma recursiva, regresa el índice
    donde se encuentra, esta función tiene dos casos base, 
    si el nodo actual es \texttt{NULL} regresa -1, si el nodo actual
    contiene el dato buscado regresa 0, en caso contrario hace una llamada
    recursiva al siguiente nodo y si el resultado es -1 regresa -1,
    si el resultado es diferente de -1 regresa el resultado + 1.

    \item \texttt{hacer\_nula(lista *l)}:
    Esta función libera la memoria de todos los nodos de la lista y 
    deja la lista vacía.

    \item \texttt{hacer\_nula\_recursivo(nodo *n)}:
    Esta función libera la memoria de todos los nodos de la lista de
    forma recursiva, tiene un caso base que es cuando el nodo actual
    es \texttt{NULL}, en este caso no hace nada y regresa, en caso 
    contrario hace una llamada recursiva al siguiente nodo y después 
    libera la memoria del nodo actual.
\end{itemize}

\section{Conclusiones}

Esta practica es una buena forma de practicar el manejo de apuntadores
y memoria dinámica en C, la implementación de algunos métodos son un 
reto que incita a pensar en el manejo correcto de los apuntadores y las 
funciones recursivas.

La implementación del cifrado Cesar aunque es sencilla permite ejercitar
el manejo de archivos y ejercitar los conocimientos básicos del lenguaje
que mas adelante se usara para implementar algoritmos mas complejos que 
impliquen el uso de concurrencia.

\end{document}